\documentclass{article}
\usepackage[utf8]{inputenc}

\title{Tarea lenguajes 2}
\author{Arteaga Hernández Alejandro Iabin}
\date{Septiembre 2017}

\usepackage{natbib}
\usepackage{graphicx}

\begin{document}

\maketitle

\section{}
\subsection{C#}
\begin{verbatim}
//'t' es de ligado
int t = 44;
//'var' es de ligado 't' es ligada
int var = t;
//'i' es de ligado 
int i = 0;
//'a' es de ligado 'something' es libre
Button a = new Button(something);
//la primera i es ligada 
for(i;i<t;i++){
    // 'y' es de ligado 'var es ligada
    int y = var;
    //'k' es libre
    Console.WriteLine(k);
}
\end{verbatim}
\subsection{Php}
\begin{verbatim}
//'k' es de ligado 'r' es libre
$k = $r;
//'$query' es de ligado
$query = "SELECT * FROM `articulo` ORDER BY fecha DESC";
//'$res' es de ligado y '$query' es ligada
$res = $this->con->query($query);
//'$index' es de ligado
$index = '';
//'$i' es de ligado
$i = 0;
//'usuario' es de ligado 
$usuario  = 4;
//'$articulos' es de ligado 
$articulos = array();
//'$res' es ligada 'row' de ligado
while($row = $res->fetch_array(MYSQLI_ASSOC)){
    //'$index' es de ligado
    $index .= "<div class='well'>";
    //'$arreglo' es de ligado '$row' es ligada
    $arreglo = explode(" ", $row['categoria_id']);
    //'nombre' es de ligado y 'usuario' ligada
    $nombre = $usuario->nombre;
    //'fechaEs' de ligado 'row' ligada
    $fechaEs = Article::fechaEs($row['fecha']);
    }
\end{verbatim}

\section{}
\begin{verbatim}
{ with { x 4}
    { with { x 5}
        { with { f { fun { y } {+ x y }}}
            { f 10}}}}
\end{verbatim}
Ben está equivocado, en este caso la evaluación estática y dinámica son las mismas (num 15), 
pero si usaramos alcance dinámico con la evaluación mas vieja, sería (num 14) lo cual es diferente.\\\\
    Esto pasa porque la evaluación dinámica es absoluta, siempre se lee desde el tope de la pila, sin importar el ambiente,mientras que
la evaluación estática, es relativa, depende del ambiente totalmente, y la disposición de la llamada, por ende, es imposible,inferir alguna de la otra, ya que las evaluaciones 
son independientes del ambiente.

\section{}
No tiene que violar la evaluación estática, como with es una función del lenguaje, puede verse de la forma 
\begin{verbatim}
    {parametros y funciones del objeto
        {parametros del with }}
\end{verbatim}
De manera que cuando se crea la pila y se evalua estáticamente, primero entran los parametros y funciones del objeto, de manera que
cuando las funciones dentro del with son ejecutadas, los parametros y funciones del with son alcanzables con alcance estático.
\section{}
Si el identificador no fue substituido en el algoritmo de substitución se le adhiere una connotación negativa, si la tiene, entonces es un identificador libre
\section{}
\subsection{}
\begin{verbatim}
    {with {x1 i}
        {with {x2 i}
            .
            .
            .
                {with {xn i}
                    {with {f {fun {y x} {+ x y}}}}
                        {f x1 {f x2 {... {f xn-1 xn}}}}}}}

 
\end{verbatim}
\subsection{}
Sustitución usando stacks,encontrar el primer elemento x1 desde f, toma n pasos, encontrar x2 toma n-1 pasos, encontrar x3 toma n-2 pasos, así hasta que encontrar xn toma 0 pasos, la suma de todos los pasos (n-1)+(n-2)+...+1+0 es de O($n^2$) como nos enseño nuestro amigo Gauss
\section{}
Supongamos que $|L|>|DL| $  sea   $ n = |L|$  \\  Como es un conjunto hay n diferentes elementos, una variable es ligada si está determinada
por una definición anterior, que le asigna un valor, esto significa que por la definición, debe haber al menos n variables de ligado que le asignan valor
a las n variables ligadas, lo cual es contradicción, por lo tanto $|L|\leq|DL|$\\\\\diamond
\bibliographystyle{plain}
\bibliography{references}

\bibitem{C++} 
FING – UDELAR  5 sept 2017
\\\texttt{https://www.fing.edu.uy/~darosa/eval08Logica-de-predicados.pdf}

\bibitem{php} 
Brown  5 sept 2017
\\\texttt{http://papl.cs.brown.edu/2014/Interpreting\_Functions.html\#\%28part.\_eager-lazy-subst\%29}

\bibitem{php} 
Etnasoft  5 sept 2017
\\\texttt{http://www.etnassoft.com/2011/06/03/estudiando-la-funcion-width-en-javascript-y\\-el-porque-no-hay-que-usarla/}


\end{document}


f 10
f {fun y {+ x y}}
x 5
x 4
